\documentclass[10pt,a4paper]{amsart}
\usepackage{amsmath,amssymb,amsthm,graphicx}
% Retained preamble for the Stacks Project chapter snapshot in this directory.
% The chapter's own preamble.tex is a copy of the shared Stacks Project preamble; the block
% below is sliced MECHANICALLY, line for line, from it. Nothing was added, removed or
% redefined; the wrapper class, the xr-hyper/hyperref wiring and the externaldocument list
% are omitted because the assembler supplies the document class.
%
% lines taken from preamble.tex: 3, 5, 9, 11, 12, 156, 157, 158, ... (51 lines total)
%
% package coverage restored from the chapter's own preamble (2026-10-13): xy, xy-2cell

\usepackage[all]{xy}
\xyoption{2cell}
\UseAllTwocells

\newenvironment{reference}{\comment}{\endcomment}
\newenvironment{slogan}{\comment}{\endcomment}
\newenvironment{history}{\comment}{\endcomment}
\theoremstyle{plain}
\newtheorem{theorem}[subsection]{Theorem}
\newtheorem{proposition}[subsection]{Proposition}
\newtheorem{lemma}[subsection]{Lemma}
\theoremstyle{definition}
\newtheorem{definition}[subsection]{Definition}
\newtheorem{example}[subsection]{Example}
\newtheorem{exercise}[subsection]{Exercise}
\newtheorem{situation}[subsection]{Situation}
\theoremstyle{remark}
\newtheorem{remark}[subsection]{Remark}
\newtheorem{remarks}[subsection]{Remarks}
\numberwithin{equation}{subsection}
\def\lim{\mathop{\mathrm{lim}}\nolimits}
\def\colim{\mathop{\mathrm{colim}}\nolimits}
\def\Spec{\mathop{\mathrm{Spec}}}
\def\Hom{\mathop{\mathrm{Hom}}\nolimits}
\def\Ext{\mathop{\mathrm{Ext}}\nolimits}
\def\SheafHom{\mathop{\mathcal{H}\!\mathit{om}}\nolimits}
\def\SheafExt{\mathop{\mathcal{E}\!\mathit{xt}}\nolimits}
\def\Sch{\mathit{Sch}}
\def\Mor{\mathop{\mathrm{Mor}}\nolimits}
\def\Ob{\mathop{\mathrm{Ob}}\nolimits}
\def\Sh{\mathop{\mathit{Sh}}\nolimits}
\def\NL{\mathop{N\!L}\nolimits}
\def\CH{\mathop{\mathrm{CH}}\nolimits}
\def\proetale{{pro\text{-}\acute{e}tale}}
\def\etale{{\acute{e}tale}}
\def\QCoh{\mathit{QCoh}}
\def\Ker{\mathop{\mathrm{Ker}}}
\def\Im{\mathop{\mathrm{Im}}}
\def\Coker{\mathop{\mathrm{Coker}}}
\def\Coim{\mathop{\mathrm{Coim}}}
\DeclareMathSymbol{\boxtimes}{\mathbin}{AMSa}{"02}
\def\QCohstack{\mathcal{QC}\!\mathit{oh}}
\def\Cohstack{\mathcal{C}\!\mathit{oh}}
\def\Spacesstack{\mathcal{S}\!\mathit{paces}}
\def\Quotfunctor{\mathrm{Quot}}
\def\Hilbfunctor{\mathrm{Hilb}}
\def\Curvesstack{\mathcal{C}\!\mathit{urves}}
\def\Polarizedstack{\mathcal{P}\!\mathit{olarized}}
\def\Complexesstack{\mathcal{C}\!\mathit{omplexes}}
\def\Pic{\mathop{\mathrm{Pic}}\nolimits}
\def\Picardstack{\mathcal{P}\!\mathit{ic}}
\def\Picardfunctor{\mathrm{Pic}}
\def\Deformationcategory{\mathcal{D}\!\mathit{ef}}
\usepackage{cleveref}
\crefname{originallemma}{original lemma}{original lemmas}
\Crefname{originallemma}{Original lemma}{Original lemmas}
\makeatletter
\newlabel{categories-definition-category-connected}{{002S}{source}{}{sourcepointer-1}{}}
\newlabel{categories-definition-category-connected@cref}{{[originallemma][0][]002S}{[1][0][]source}{}{}{}}
\makeatother
\crefname{originallemma}{original lemma}{original lemmas}
\Crefname{originallemma}{Original lemma}{Original lemmas}
\makeatletter
\newlabel{categories-lemma-connected-limit-over-X}{{002T}{source}{}{sourcepointer-2}{}}
\newlabel{categories-lemma-connected-limit-over-X@cref}{{[originallemma][0][]002T}{[1][0][]source}{}{}{}}
\makeatother
\crefname{originallemma}{original lemma}{original lemmas}
\Crefname{originallemma}{Original lemma}{Original lemmas}
\makeatletter
\newlabel{definition-covering-SR}{{01G3}{source}{}{sourcepointer-3}{}}
\newlabel{definition-covering-SR@cref}{{[originallemma][0][]01G3}{[1][0][]source}{}{}{}}
\makeatother
\crefname{originallemma}{original lemma}{original lemmas}
\Crefname{originallemma}{Original lemma}{Original lemmas}
\makeatletter
\newlabel{lemma-coprod-prod-SR}{{01G2}{source}{}{sourcepointer-4}{}}
\newlabel{lemma-coprod-prod-SR@cref}{{[originallemma][0][]01G2}{[1][0][]source}{}{}{}}
\makeatother
\crefname{originallemma}{original lemma}{original lemmas}
\Crefname{originallemma}{Original lemma}{Original lemmas}
\makeatletter
\newlabel{simplicial-example-cosk0}{{0182}{source}{}{sourcepointer-5}{}}
\newlabel{simplicial-example-cosk0@cref}{{[originallemma][0][]0182}{[1][0][]source}{}{}{}}
\makeatother
\crefname{originallemma}{original lemma}{original lemmas}
\Crefname{originallemma}{Original lemma}{Original lemmas}
\makeatletter
\newlabel{simplicial-lemma-existence-cosk}{{0183}{source}{}{sourcepointer-6}{}}
\newlabel{simplicial-lemma-existence-cosk@cref}{{[originallemma][0][]0183}{[1][0][]source}{}{}{}}
\makeatother
\crefname{originallemma}{original lemma}{original lemmas}
\Crefname{originallemma}{Original lemma}{Original lemmas}
\makeatletter
\newlabel{simplicial-section-skeleton}{{017Z}{source}{}{sourcepointer-7}{}}
\newlabel{simplicial-section-skeleton@cref}{{[originallemma][0][]017Z}{[1][0][]source}{}{}{}}
\makeatother
\crefname{originallemma}{original lemma}{original lemmas}
\Crefname{originallemma}{Original lemma}{Original lemmas}
\makeatletter
\newlabel{sites-lemma-covering-surjective-after-sheafification}{{00WT}{source}{}{sourcepointer-8}{}}
\newlabel{sites-lemma-covering-surjective-after-sheafification@cref}{{[originallemma][0][]00WT}{[1][0][]source}{}{}{}}
\makeatother
\title{Sheafification turns hypercovering conditions into surjections of F(K)}
\author{The Stacks Project Authors}
\date{stacks-project:hypercovering.tex:local-snapshot-20261012; GNU FDL 1.2 or later}
\begin{document}
\maketitle
\noindent\textit{Editorial note.} Statement span is lines 684-708 of sources/stacks-hypercovering/pinned-original-source.tex; proof span is lines 710-760, the complete author proof (49 body lines).\par\medskip
\noindent\textit{Editorial note.} Navigation stubs: 8 cited-but-not-included labels are supplied as non-load-bearing pointers with their pinned tags where the snapshot tag index has them.\par\medskip
\noindent\textit{Editorial note.} extra\_packages is empty on purpose: the retained preamble already loads xy with options and a second declaration clashes in the assembler wrapper.\par\medskip
\noindent\textit{Editorial note.} Text is retained verbatim, including the author's own typos and placeholders.\par\medskip
\noindent\textit{Editorial note.} Queue index 8 of queue\_hypercovering.json (pb=2035, stmt 684-708, proof 710-760).\par\medskip
\noindent\textit{Editorial note.} stacks-hypercovering wave-1 (hc01): fresh adjudication accept (2026-10-14); pool per standing fallback order (hypercovering -> pic -> spaces-more-cohomology) after md01 closed the moduli pool; hypercovering.tex 53 labels/34 with proof judged 26 accept / 8 reject / 14 pb=0 env-dup rows excluded; spans cite-free and reference/slogan/history-free so no patchenv needed. Named omission: Connectedness of (Delta/[m])\_\{<=n\} for n >= 1 is left as an 'amusing exercise' - a named small fact; the F'/finite-limits argument is in-span.\par\medskip
\noindent\textit{Allowed standard prerequisites (editorial).}\begin{itemize}
\item Explicit prerequisite boundary. The cited inputs are prerequisites; the argument in the retained span is not replaced by any cited result.
\end{itemize}\medskip
\ifdefined\hypertarget\hypertarget{sourcepointer-1}{}\fi
\noindent\textit{Source reference.} Original lemma 002S: sources/stacks-hypercovering/pinned-original-source.tex, source lines 760--760. Cited by the retained statement or proof but not part of the collected spans; supplied only as a navigation pointer. The argument does not depend on its content.\par\medskip
\ifdefined\hypertarget\hypertarget{sourcepointer-2}{}\fi
\noindent\textit{Source reference.} Original lemma 002T: sources/stacks-hypercovering/pinned-original-source.tex, source lines 760--760. Cited by the retained statement or proof but not part of the collected spans; supplied only as a navigation pointer. The argument does not depend on its content.\par\medskip
\ifdefined\hypertarget\hypertarget{sourcepointer-3}{}\fi
\noindent\textit{Source reference.} Original lemma 01G3: sources/stacks-hypercovering/pinned-original-source.tex, source lines 510--521. Cited by the retained statement or proof but not part of the collected spans; supplied only as a navigation pointer. The argument does not depend on its content.\par\medskip
\ifdefined\hypertarget\hypertarget{sourcepointer-4}{}\fi
\noindent\textit{Source reference.} Original lemma 01G2: sources/stacks-hypercovering/pinned-original-source.tex, source lines 391--404. Cited by the retained statement or proof but not part of the collected spans; supplied only as a navigation pointer. The argument does not depend on its content.\par\medskip
\ifdefined\hypertarget\hypertarget{sourcepointer-5}{}\fi
\noindent\textit{Source reference.} Original lemma 0182: sources/stacks-hypercovering/pinned-original-source.tex, source lines 760--760. Cited by the retained statement or proof but not part of the collected spans; supplied only as a navigation pointer. The argument does not depend on its content.\par\medskip
\ifdefined\hypertarget\hypertarget{sourcepointer-6}{}\fi
\noindent\textit{Source reference.} Original lemma 0183: sources/stacks-hypercovering/pinned-original-source.tex, source lines 760--760. Cited by the retained statement or proof but not part of the collected spans; supplied only as a navigation pointer. The argument does not depend on its content.\par\medskip
\ifdefined\hypertarget\hypertarget{sourcepointer-7}{}\fi
\noindent\textit{Source reference.} Original lemma 017Z: sources/stacks-hypercovering/pinned-original-source.tex, source lines 760--760. Cited by the retained statement or proof but not part of the collected spans; supplied only as a navigation pointer. The argument does not depend on its content.\par\medskip
\ifdefined\hypertarget\hypertarget{sourcepointer-8}{}\fi
\noindent\textit{Source reference.} Original lemma 00WT: sources/stacks-hypercovering/pinned-original-source.tex, source lines 760--760. Cited by the retained statement or proof but not part of the collected spans; supplied only as a navigation pointer. The argument does not depend on its content.\par\medskip
\begin{lemma}
\label{lemma-hypercovering-F}
Let $\mathcal{C}$ be a site with fibre products.
Let $X \in \Ob(\mathcal{C})$ be an object of $\mathcal{C}$.
Let $K$ be a hypercovering of $X$.
Consider the simplicial object $F(K)$ of $\textit{PSh}(\mathcal{C})$,
endowed with its augmentation to the constant simplicial presheaf $h_X$.
\begin{enumerate}
\item The morphism of presheaves $F(K)_0 \to h_X$ becomes
a surjection after sheafification.
\item The morphism
$$
(d^1_0, d^1_1) :
F(K)_1
\longrightarrow
F(K)_0 \times_{h_X} F(K)_0
$$
becomes a surjection after sheafification.
\item For every $n \geq 1$ the morphism
$$
F(K)_{n + 1} \longrightarrow (\text{cosk}_n \text{sk}_n F(K))_{n + 1}
$$
turns into a surjection after sheafification.
\end{enumerate}
\end{lemma}

\begin{proof}
We will use the fact that if
$\{U_i \to U\}_{i \in I}$ is a covering of the site
$\mathcal{C}$, then the morphism
$$
\amalg_{i \in I} h_{U_i} \to h_U
$$
becomes surjective after sheafification, see
Sites, Lemma \ref{sites-lemma-covering-surjective-after-sheafification}.
Thus the first assertion follows immediately.

\medskip\noindent
For the second assertion, note that according to
Simplicial, Example \ref{simplicial-example-cosk0}
the simplicial object $\text{cosk}_0 \text{sk}_0 K$
has terms $K_0 \times \ldots \times K_0$. Thus
according to the definition of a hypercovering we
see that $(d^1_0, d^1_1) : K_1 \to K_0 \times K_0$ is a
covering. Hence (2) follows from the claim above
and the fact that $F$ transforms products into fibred
products over $h_X$.

\medskip\noindent
For the third, we claim that
$\text{cosk}_n \text{sk}_n F(K) =
F(\text{cosk}_n \text{sk}_n K)$ for $n \geq 1$.
To prove this, denote temporarily $F'$ the functor
$\text{SR}(\mathcal{C}, X) \to \textit{PSh}(\mathcal{C})/h_X$.
By Lemma \ref{lemma-coprod-prod-SR} the functor
$F'$ commutes with finite limits.
By our description of the $\text{cosk}_n$ functor in
Simplicial, Section \ref{simplicial-section-skeleton}
we see that $\text{cosk}_n \text{sk}_n F'(K) =
F'(\text{cosk}_n \text{sk}_n K)$.
Recall that the category used in the description of
$(\text{cosk}_n U)_m$ in
Simplicial, Lemma \ref{simplicial-lemma-existence-cosk}
is the category $(\Delta/[m])^{opp}_{\leq n}$. It is an
amusing exercise to show that $(\Delta/[m])_{\leq n}$ is
a connected category (see
Categories, Definition \ref{categories-definition-category-connected})
as soon as $n \geq 1$. Hence,
Categories, Lemma \ref{categories-lemma-connected-limit-over-X}
shows that $\text{cosk}_n \text{sk}_n F'(K) =
\text{cosk}_n \text{sk}_n F(K)$. Whence the claim.
Property (2) follows from this, because now we see that
the morphism in (2) is the result of applying the
functor $F$ to a covering as in Definition \ref{definition-covering-SR},
and the result follows from the first fact mentioned
in this proof.
\end{proof}

\end{document}
